\subsection{Audio Generation}\label{sec:related_audio}

\textbf{Video-to-audio generation.}
There are many recent studies exploring video-to-audio generation. Most of them are based on latent diffusion models similar to ours~\citep{luo2024diff, xu2024video, xing2024seeing, zhang2024foleycrafter} with a few exceptions being token-based language models~\citep{kondratyuk2023videopoet, mei2023foleygen}. Diff-Foley~\citep{luo2024diff} and VTA-LDM~\citep{xu2024video} are the standard latent diffusion models with U-Net architectures conditioned on video features, extracted from a pre-trained contrastive audio-video encoder (CAVP~\citep{luo2024diff}) and video-text encoder (CLIP4CLIP~\citep{luo2022clip4clip}), respectively. Seeing-and-hearing (S\&H)~\citep{xing2024seeing} proposes a training-free method that uses ImageBind as a classifier guidance to guide a pre-trained diffusion-based text-to-audio (TTA) model (AudioLDM~\citep{liu2023audioldm}) to generate video aligned audio. FoleyCrafter~\citep{zhang2024foleycrafter} uses the adaptor-based approach~\citep{zhang2023adding} to finetune a pre-trained TTA model to add video control. To enhance the temporal alignment between audio and video, it further conditions on timestamps to inform the model which segments are sounded/silence, which are predicted from the video during inference.

Most of these models are directly trained on, or built on TTA models trained on solely in-the-wild datasets, such as VGGSound (550 hours)~\citep{vggsound} or AudioSet (5K hours)~\citep{gemmeke2017audio}.
These datasets have several limitations. In terms of quality, many of the videos are recorded with non-professional devices like smartphones or low-end cameras, and hence both the video and the audio quality are subpar.
In terms of sound design, most of the videos are uploaded by amateur creators who had done none or minimal post-processing and post-production, which contain only diegetic sounds. Compared to professionally created films, such videos contain an abundance of distracting sounds (\eg, irrelevant off-screen speech, wind noise, high ambient noise), do not emphasize main sound events (\eg, exaggerated breathing sound, Foley sounds for footsteps and object clatters), and lack carefully designed non-diegetic sounds that are critical for a cinematic feeling (\eg, use of riser, braam, underscore music, action scoring).
Training on such datasets inevitably prevents the resulting model from generating cinematic soundtracks including both music and sound effects.

On the other hand, the size of prior video-to-audio generation models are relatively small, typically ranging from 300M to 1.3B parameters~\citep{xing2024seeing, luo2024diff, mei2023foleygen}. Combined with the data size, the scale also limits the performance of these models, as we demonstrated in the ablation study (see~\cref{sec:audio_ablation}) that scaling to 13B significantly improves both quality and video alignment.
Compared with the prior works that also offer text control, we additionally provide quality control and fine-grained music control, which alleviates the quality issue when training on mixed-quality data, and improves soundtrack design flexibility.

There are a few products offering video-to-audio capabilities, including PikaLabs\footnote{\url{https://pika.art/}} and ElevenLabs.\footnote{\url{https://github.com/elevenlabs/elevenlabs-examples/tree/main/examples/sound-effects/video-to-sfx}}, but neither can really generate motion-aligned sound effects or cinematic soundtracks with both music and sound effects. PikaLabs supports sound effect generation with video and optionally text prompts; however it will generate audio longer than the video where a user needs to select an audio segment to use. This implies under the hood it may be an audio generation model conditioned on a fixed number of key image frames. The maximum audio length is capped at 15 seconds without joint music generation and audio extension capabilities, preventing its application to soundtrack creation for long-form videos.
ElevenLabs leverages GPT-4o to create a sound prompt given four image frames extracted from the video (one second apart), and then generates audio using a TTA model with that prompt.
Lastly, Google released a research blog\footnote{\url{https://deepmind.google/discover/blog/generating-audio-for-video/}} describing their video-to-audio generation models that also provide text control. Based on the video samples, the model is capable of sound effects, speech, and music generation. However, the details (model size, training data characterization) about the model and the number of samples (13 samples with 11 distinct videos) are very limited, and no API is provided. It is difficult to conclude further details other than the model is diffusion-based and that the maximum audio length may be limited as the longest sample showcased is less than 15 seconds.

\textbf{Video to music generation.}
Many studies often focus on symbolic music (MIDI) generation~\citep{Di2021VideoBM, Zhuo2022VideoBM, Kang2023Video2MusicSM} for piano or other instruments, as MIDI is easier to predict compared to raw audio. Compared to end-to-end modeling, such a paradigm imposes many restrictions. First, MIDI is a form of music transcription which cannot capture all the details from the original music. Hence, the generated music from such systems tend to sound more monotonic. Second, it requires MIDI annotation or a high quality music transcription model, which limits the sources of training data one can consider. Lastly, such models cannot learn the relationship between music and other audio components like sound effects and speech, which are not trivial in cinematic films.

Most of these works, along those directly predicting audio~\citep{Zhu2022QuantizedGF}, extract low-level music-related features from videos, such as human motion, scene change timing, and tempo, for conditioning. These features along with the training data (\eg, dancing videos) are rather domain specific, which cannot generalize to general videos.

Our work is most related to \citet{Su2023V2MeowMT} and \citet{Tian2024VidMuseAS}. Both prior works extract general video features (CLIP, flow, image tokens) and predict general audio representations (EnCodec~\citep{Defossez2022HighFN}, Soundstream~\citep{Zeghidour2022SoundStreamAE}, w2vBERT~\citep{Chung2021w2vBERTCC}). In addition to our novel scaling, the main differences in our work are twofold: first we aim for joint non-diegetic music and sound effect generation while these studies only focus on non-diegetic music; second we adopt diffusion modeling which is free of tokenization information loss and hence enjoys the other benefits described in previous sections, while \citet{Tian2024VidMuseAS} points out explicitly that their quality is suffer from the audio codec limitation.
